\honest{Fake Explanations}{Eliezer Yudkowsky}{2007}

I start with a joke. An instructor shows physics students a metal plate next to a radiator. The side near the radiator feels cool and the far side warm. The students suggest convection, strange metals and other causes; none says ``I don't know'' or ``This seems impossible.'' The instructor had turned the plate around.

Take the student who stammers, ``Eh, maybe because of the heat conduction and so?'' The words are easy to say loudly; do they control what the student expects? Put ``phlogiston'' or ``magic'' after ``because of'' instead. ``Heat conduction'' is what Spock would say, and ``magic'' what Giles would say; the phrases belong to different ``literary genres.'' A Bayesian ignores genre: what matters in a model is how it constrains what you expect. \nb{Richard Feynman made the same comparison. In \textsc{Surely You're Joking, Mr. Feynman!} (1985) he wrote of a textbook's ``Energy makes it go'': ``Suppose it's `Wakalixes.' That's the general principle: `Wakalixes makes it go.' There's no knowledge coming in.'' The post does not mention him.}

Heat conduction normally leads you to expect the near side to be warmer. If it can also explain the near side being cooler, it can explain ``pretty much anything.'' \nb{With the right starting condition it does explain a cooler near side, and the post's own physicist, below, explains the plate that way. What the student's answer lacks is the starting condition.} An answer that fits any outcome carries ``zero knowledge''; I call ``heat conduction,'' used this way, ``a disguised hypothesis of maximum entropy.'' \nb{The principle is Karl Popper's.}

Instead of guessing, one could measure the plate and use the diffusion equation. I am ``not physicist enough'' to say how many measurements that would take. A physicist would measure and say that the plate was in equilibrium two and a half minutes ago, was turned around, and is ``now approaching equilibrium again.''

The students' deeper error was to think they were doing physics. ``They simply moved their magic from one literary genre to another.''
